% !TEX root = ../Report.tex
% ======================================================================
\chapter{Methodology}
\label{ch:methodology}
% ======================================================================

This chapter explains the tool selection, testing
methodology, and evaluation design used to deliver GiTrip. The development process and methodology are discussed in the Project Management Report.

\section{Choice of Tools and Technologies}
GiTrip's major technology choices were guided by the requirement to deliver a
complete, deployable web system with both server-rendered UI and asynchronous
planning/merge operations.

\subsection{Node.js and Express}
GiTrip uses a JavaScript full-stack to reduce development overhead in a
single-developer project and share data models and validation logic across
client/server. Node.js'~\cite{ref_node_about} non-blocking I/O suits the I/O-heavy workload (routing/geocoding calls plus repository
read/write operations). Express provides a minimal server and API layer since EJS supports rapid iteration on key UX
flows without a complex front-end build chain~\cite{ref_express}. This combination
enables server-rendered views (for fast initial load and discoverability) while
still allowing interactive enhancements via client-side JavaScript.

\subsection{SQLite for Persistent Storage}
SQLite was selected as a lightweight relational store with transactional integrity and low operational complexity~\cite{ref_sqlite_about}. For
the expected scale (tens of stops per trip; snapshots typically $<100$\,KB),
snapshot storage is practical while simplifying branching, rollback, and
merging. Its file-based
deployment model aligns with container deployment and makes backups simple. GiTrip's collaboration model (branch-and-merge) does not require
multi-writer database replication; concurrency is managed at the application
level.

\subsection{Mapping and Geocoding}
Place search uses Nominatim~\cite{ref_nominatim}, and Leaflet~\cite{ref_leaflet} with OpenStreetMap~\cite{ref_osm} provides interactive mapping.
Photon~\cite{ref_photon} (komoot) was added as a fuzzy-search fallback: when Nominatim returns no
results for a query, Photon's Elasticsearch-based index over the same OpenStreetMap (OSM) data
handles typos and punctuation variants that Nominatim's exact matching misses.
A server-side proxy route is used to apply rate limiting and caching to avoid
overloading the upstream services.

\subsection{Routing Providers}
Travel-time computation uses Openrouteservice (ORS) matrix endpoints~\cite{ref_ors_matrix} (for matrix-enabled modes) and Google Directions~\cite{ref_google_directions_api} (for transit). A key design requirement was provider fallback:
when API calls fail or quota limits are reached, GiTrip degrades gracefully
using Haversine-based estimates with mode-dependent speed assumptions and
constant-gap defaults. API quota outage is mitigated via caching, throttling, and distance-based fallback estimates.

\subsection{Usability Frameworks}
UI design and evaluation were guided by Nielsen's usability heuristics, WCAG 2.2, SUS, and TAM-inspired constructs, as discussed in Chapter~\ref{ch:background}.

\section{Testing Strategy}
Testing was treated as a first-class requirement due to the complexity of merge
logic and scheduling constraints. The strategy combined:
\begin{itemize}
  \item \textbf{Unit tests} for opening-hours parsing, time utilities, and merge
        conflict detection.
  \item \textbf{Functional tests} covering end-to-end workflows (plan, commit,
        branch, merge, resolve, export).
  \item \textbf{Non-functional checks} covering security headers, rate limiting,
        and accessibility considerations.
\end{itemize}

All unit tests use the Node.js native test framework (\texttt{node:test}) with
strict assertions.

\section{Evaluation Design}
Evaluation was conducted via:
\begin{itemize}
  \item \textbf{User Acceptance Testing (UAT)} on Evaluation Day with an observer
        checklist recording objective task completion, time-to-finish, and hints
        given.
  \item \textbf{Post-demo survey} on Evaluation Day, collecting Likert responses
        (including SUS), Git familiarity, and free-text feedback.
\end{itemize}

This design intentionally combines objective evidence (task completion and time)
with subjective user perception (usability and perceived value), aligning with a
pragmatic interpretation of TAM in small-sample contexts.

\section{Ethics and Data Handling}
The evaluation collected minimal personal data. Survey responses were anonymous,
and the system does not require real personal identity information for most
features. Where authentication is used, sessions are stored via random
identifiers and protected by HttpOnly cookies. GDPR~\cite{ref_gdpr} considerations are discussed
in Chapter~\ref{ch:design}.
